Con la distancia base de Chocontá a Tunja (17.69 km, valor oficial registrado de un sector truncado), el desvío por Santander ahorra 70\,600 COP de peaje y suma 345.15 km. La ruta por Bogotá es más barata en dinero si el costo variable por km supera el ahorro dividido entre los km extra, 204.5 COP/km. Con la cota inferior (56.96 km, la línea recta entre los nodos) el umbral es 230.8 COP/km y con las progresivas (61.0 km, con un hueco de 49.1 km sin sector registrado y el extremo a 11 km de Tunja), 233.9 COP/km. Con el costo variable vigente, 1\,157.0 COP/km (combustible con el precio de la CREG y el consumo de la UPME, insumos verificados), la ruta por Bogotá es la más barata en dinero en los tres casos, y lo supera 5.7, 5.0, 4.9 veces, respectivamente. Los otros costos por km (llantas, lubricantes, mantenimiento y conductor) no tienen fuente y se toman iguales a 0 (supuesto del equipo), de modo que el costo variable real no es menor que el calculado. No se incluyen tiempo ni riesgo.
